\section{Pre-training}

Our language model pre-training process consists of several key components. First, we carefully curate high-quality training data through sophisticated filtering and scoring mechanisms, combined with strategic data mixture. Second, we conduct extensive research on hyperparameter optimization to effectively train models at various scales. Finally, we incorporate specialized long-context pre-training to enhance the model's ability to process and understand extended sequences. Below, we detail our approaches to data preparation, hyperparameter selection, and long-context training.

\label{sec:pre}

\subsection{Pre-training Data}

Qwen2.5 demonstrates significant enhancements in pre-training data quality compared to its predecessor Qwen2. These improvements stem from several key aspects:

\begin{enumerate}[label=(\arabic*)]
    \item \textbf{Better data filtering}. High-quality pre-training data is crucial for model performance, making data quality assessment and filtering a critical component of our pipeline. We leverage Qwen2-Instruct models as data quality filters that perform comprehensive, multi-dimensional analysis to evaluate and score training samples. The filtering method represents a significant advancement over our previous approach used for Qwen2, as it benefits from Qwen2's expanded pre-training on a larger multilingual corpus. The enhanced capabilities enable more nuanced quality assessment, resulting in both improved retention of high-quality training data and more effective filtering of low-quality samples across multiple languages.

    \item \textbf{Better math and code data}. During the pre-training phase of Qwen2.5, we incorporate training data from Qwen2.5-Math~\citep{qwen2.5math} and Qwen2.5-Coder~\citep{qwen2.5coder}. This data integration strategy proves highly effective, as these specialized datasets are instrumental in achieving state-of-the-art performance on mathematical and coding tasks. By leveraging these high-quality domain-specific datasets during pre-training, Qwen2.5 inherits strong capabilities in both mathematical reasoning and code generation.

    \item \textbf{Better synthetic data}. To generate high-quality synthetic data, particularly in mathematics, code, and knowledge domains, we leverage both Qwen2-72B-Instruct~\citep{qwen2} and Qwen2-Math-72B-Instruct~\citep{qwen2math}. The quality of this synthesized data is further enhanced through rigorous filtering using our proprietary general reward model and the specialized Qwen2-Math-RM-72B~\citep{qwen2math} model.

    \item \textbf{Better data mixture}. To optimize the pre-training data distribution, we employ Qwen2-Instruct models to classify and balance content across different domains. Our analysis revealed that domains like e-commerce, social media, and entertainment are significantly overrepresented in web-scale data, often containing repetitive, template-based, or machine-generated content. Conversely, domains such as technology, science, and  academic research, while containing higher-quality information, are traditionally underrepresented. Through strategic down-sampling of overrepresented domains and up-sampling of high-value domains, we ensure a more balanced and information-rich training dataset that better serves our model's learning objectives.
\end{enumerate}

Building on these techniques, we have developed a larger and higher-quality pre-training dataset, expanding from the 7 trillion tokens used in Qwen2~\citep{qwen2} to \textbf{18 trillion} tokens.

\subsection{Scaling Law for Hyper-parameters}
We develop scaling laws for hyper-parameter based on the pre-training data of Qwen2.5~\citep{chinchilla,kaplanscaling}. While previous studies~\citep{llama3,falcon,chinchilla} primarily used scaling laws to determine optimal model sizes given compute budgets, we leverage them to identify optimal hyperparameters across model architectures. Specifically, our scaling laws help determine key training parameters like batch size $B$ and learning rate $\mu$ for both dense models and MoE models of varying sizes.

Through extensive experimentation, we systematically study the relationship between model architecture and optimal training hyper-parameters. Specifically, we analyze how the optimal learning rate $\mu_\textrm{opt}$ and batch size $B_\textrm{opt}$ vary with model size $N$ and pre-training data size $D$. Our experiments cover a comprehensive range of architectures, including dense models with 44M to 14B parameters and MoE models with 44M to 1B activated parameters, trained on datasets ranging from 0.8B to 600B tokens. Using these optimal hyper-parameter predictions, we then model the final loss as a function of model architecture and training data scale.

Additionally, we leverage scaling laws to predict and compare the performance of MoE models with varying parameter counts against their dense counterparts. This analysis guides our hyper-parameter configuration for MoE models, enabling us to achieve performance parity with specific dense model variants (such as Qwen2.5-72B and Qwen2.5-14B) through careful tuning of both activated and total parameters.



\subsection{Long-context Pre-training}

For optimal training efficiency, Qwen2.5 employs a two-phase pre-training approach: an initial phase with a 4,096-token context length, followed by an extension phase for longer sequences. Following the strategy used in Qwen2, we extend the context length from 4,096 to 32,768 tokens during the final pre-training stage for all model variants except Qwen2.5-Turbo. Concurrently, we increase the base frequency of RoPEfrom 10,000 to 1,000,000 using the ABF technique~\citep{ropeabf}.

For Qwen2.5-Turbo, we implement a progressive context length expansion strategy during training, advancing through four stages: 32,768 tokens, 65,536 tokens, 131,072 tokens, and ultimately 262,144 tokens, with a RoPE base frequency of 10,000,000. At each stage, we carefully curate the training data to include 40\% sequences at the current maximum length and 60\% shorter sequences. This progressive training methodology enables smooth adaptation to increasing context lengths while maintaining the model's ability to effectively process and generalize across sequences of varying lengths.


To enhance our models' ability to process longer sequences during inference, we implement two key strategies: YARN~\citep{yarn} and Dual Chunk Attention~(DCA, \citealp{chunkllama}). Through these innovations, we achieve a four-fold increase in sequence length capacity, enabling Qwen2.5-Turbo to handle up to \textbf{1 million} tokens and other models to process up to 131,072 tokens. Notably, these approaches not only improve the modeling of long sequences by reducing perplexity but also maintain the models' strong performance on shorter sequences, ensuring consistent quality across varying input lengths.
